\section{Perturbing a finite-order Hamiltonian map}
\label{sec:reduction}

A short Hamiltonian perturbation of a finite-order map can replace its
fixed submanifold by the critical set of a function on that submanifold.
The tangent-space hypothesis below is the usual \emph{clean fixed-set}
condition; we state it explicitly to identify the geometric input.

\begin{proposition}
\label{prop:reduction}
Let \((M,\omega)\) be a closed symplectic manifold, and let
\(A\in\Ham(M,\omega)\) satisfy \(A^m=\id\) for some positive integer
\(m\). Suppose that \(F=\Fix(A)\) is a closed embedded submanifold and
that
\[
  T_qF=\ker(dA_q-\id)
  \qquad\text{for every }q\in F.
\]
For every smooth function \(f:F\to\R\), there is a smooth
\(A\)-invariant extension \(K:M\to\R\) such that, if \(B_s\) is its
Hamiltonian flow, then for every sufficiently small \(\varepsilon>0\)
the map
\[
  \phi_\varepsilon=A\circ B_\varepsilon
\]
is Hamiltonian and satisfies
\[
  \Fix(\phi_\varepsilon)=\CritSet(f),
\]
where the critical set on the right is viewed as a subset of \(M\).
\end{proposition}

The proof combines finite-group averaging with the exclusion of short
nonconstant periodic trajectories. We establish the latter fact first.
Lipschitz control of a vector field gives a lower bound for its
nonconstant periods, a classical principle for which Yorke proved the
sharp Euclidean estimate~\cite{Yorke1969}. The following elementary
local argument suffices here.

\begin{lemma}
\label{lem:short-periods}
Let \(X\) be a smooth vector field on a closed manifold \(M\), and let
\(\Psi_t\) be its flow. There is a number \(\delta>0\) such that
\[
  \Fix(\Psi_t)=\{q\in M:X(q)=0\}
  \qquad\text{for every }0<t<\delta.
\]
\end{lemma}

\begin{proof}
The flow is defined for all time because \(M\) is compact. Suppose that
there were points \(q_j\) and times \(t_j>0\) with
\[
  t_j\longrightarrow0,\qquad
  \Psi_{t_j}(q_j)=q_j,\qquad X(q_j)\ne0.
\]
After passing to a subsequence, assume that \(q_j\to q\). Choose
coordinates near \(q\), and a Euclidean ball about its coordinate image
whose closure lies inside the coordinate domain. By continuity of the
flow, all sufficiently late trajectory segments
\(\{\Psi_s(q_j):0\le s\le t_j\}\) lie in this ball. The coordinate vector
field \(V\) is Lipschitz there with a constant \(L\ge0\), since its
derivative is bounded on the closed ball.

Write \(u_j(s)\) for the coordinate trajectory and set
\[
  C_j=\max_{0\le s\le t_j}|\dot u_j(s)|>0.
\]
For any \(s,r\in[0,t_j]\), convexity of the ball and the Lipschitz
estimate give
\[
 \begin{split}
  |\dot u_j(s)-\dot u_j(r)|
   &=|V(u_j(s))-V(u_j(r))|\\
   &\le L|u_j(s)-u_j(r)|
    \le LC_j|s-r|
    \le LC_jt_j.
 \end{split}
\]
The closed trajectory has zero mean coordinate velocity:
\(\int_0^{t_j}\dot u_j(r)\,dr=0\). At a time \(s_j\) of maximum speed,
it follows that
\[
  C_j
  =\left|
     \frac1{t_j}\int_0^{t_j}
        \bigl(\dot u_j(s_j)-\dot u_j(r)\bigr)\,dr
    \right|
  \le LC_jt_j.
\]
Dividing by \(C_j\) contradicts \(t_j\to0\). Thus all fixed points of a
sufficiently short positive-time map are zeros of \(X\). Conversely,
every zero of \(X\) is stationary and hence fixed by every \(\Psi_t\).
\end{proof}

\begin{proof}[Proof of Proposition~\ref{prop:reduction}]
The first step is to extend $f$ so that criticality along $F$ is equivalent
to ambient criticality at points of $F$. This is the finite-group case
of symmetric criticality~\cite{Palais1979}, which we verify directly.
Extend \(f\) smoothly to a function \(\widetilde f\) on \(M\): in
submanifold charts, extend independently of the normal coordinates,
and combine these local extensions using a partition of unity. Define
\[
  K=\frac1m\sum_{j=0}^{m-1}\widetilde f\circ A^j.
\]
Then \(K\circ A=K\) and \(K|_F=f\).

Fix \(q\in F\), put \(D=dA_q\), and consider
\[
  P_q=\frac1m\sum_{j=0}^{m-1}D^j.
\]
Since \(D^m=\id\), we have \(DP_q=P_q\). Moreover, \(P_qv=v\) whenever
\(Dv=v\). Thus \(P_q\) is a projection onto
\(\ker(D-\id)=T_qF\). Differentiating the averaged extension gives,
for every \(v\in T_qM\),
\[
  dK_q(v)=d\widetilde f_q(P_qv)=df_q(P_qv).
\]
Consequently,
\begin{equation}
\label{eq:critical-restriction}
  \CritSet(K)\cap F=\CritSet(f).
\end{equation}

Because \(A\) is symplectic and \(K\circ A=K\), its differential carries
the Hamiltonian vector field \(X_K\) to itself. Uniqueness of solutions
therefore gives \(A\circ B_s=B_s\circ A\) for every \(s\). Apply
Lemma~\ref{lem:short-periods} to \(X_K\), and choose
\(0<\varepsilon<\delta/m\). We then have
\[
  \phi_\varepsilon^m=B_{m\varepsilon},
  \qquad
  \Fix(B_{m\varepsilon})=\CritSet(K).
\]
If \(\phi_\varepsilon(q)=q\), then \(q\) is fixed by
\(B_{m\varepsilon}\), so \(dK_q=0\). It follows that
\(B_\varepsilon(q)=q\), and hence \(A(q)=q\). Conversely, every point
of \(\CritSet(K)\cap F\) is fixed by \(\phi_\varepsilon\).
Equation~\ref{eq:critical-restriction} proves the asserted equality
of fixed and critical sets.

For completeness, the composition is generated by a smooth
Hamiltonian on the unit time interval. Choose a Hamiltonian path
\(A_t\) from \(\id\) to \(A\), generated by \(H_t\), and a smooth
nondecreasing function \(a:[0,1]\to[0,1]\) equal to \(0\) near \(0\)
and to \(1\) near \(1\). The time-dependent Hamiltonian
\[
  L_t(x)=
  \begin{cases}
    2\varepsilon a'(2t)K(x),&0\le t\le\tfrac12,\\[2pt]
    2a'(2t-1)H_{a(2t-1)}(x),&\tfrac12\le t\le1
  \end{cases}
\]
is smooth, since both pieces vanish near the joining time. Its flow
first follows \(B_{\varepsilon a(2t)}\), and then
\(A_{a(2t-1)}\circ B_\varepsilon\), so its time-one map is
\(\phi_\varepsilon\).
\end{proof}

The proposition imposes no bound on the total number of critical
points of \(K\), and requires invariance under $A$, not under an isotopy
from $\id$ to $A$. A critical point \(q\notin F\) satisfies
\(\phi_\varepsilon(q)=A(q)\ne q\), so it contributes no fixed point.
